\chapter{BV representatives and perimeter on an interval}\label{sec:bv-interval-structure}
The measure definition of \cite{Sim18}, Chapter 2, Section 2, becomes particularly explicit on a bounded interval $I=(a,b)$. The proofs below develop this one-dimensional specialization directly. All endpoint traces are one-sided limits of the specified representative; arbitrary values assigned at single points have no effect on $Du$.

\begin{lemma}[The cumulative function of a signed measure]\label{lem:interval-measure-primitive}
Let $\mu$ be a finite signed Radon measure on $(a,b)$. Define $F(x)=\mu((a,x])$. Then $F$ is bounded and right-continuous, has left limit $F(x-)=\mu((a,x))$, and satisfies $DF=\mu$ in distributions. Its endpoint limits are $F(a+)=0$ and $F(b-)=\mu((a,b))$.
\end{lemma}
\begin{proof}
The bound $|F|\leq|\mu|(I)$ is immediate. Continuity from above and below for the finite positive measure $|\mu|$ gives the asserted one-sided limits for $\mu$. For $\phi\in C_c^1(I)$, Fubini applies to the finite product of $|\mu|$ and Lebesgue measure and gives
\[
 -\int_a^bF(x)\phi'(x)\,dx
 =-\int_I\left(\int_t^b\phi'(x)\,dx\right)d\mu(t)
 =\int_I\phi(t)\,d\mu(t).
\]
This proves the distributional identity with the prescribed sign.
\end{proof}

\begin{theorem}[The regulated representative of an interval BV function]\label{thm:interval-bv-representative}
For $u\in BV(I)$, with $\mu=Du$, there is a constant $c$ such that the right-continuous representative
\[
 u^+(x)=c+\mu((a,x])
\]
equals $u$ almost everywhere. It has left limits $u^-(x)=c+\mu((a,x))$, finite endpoint traces, and
\[
 u^+(x)-u^-(x)=\mu(\{x\}).
\]
There are at most countably many nonzero jumps. Its pointwise variation, defined by the supremum of $\sum_{j=1}^N|u^+(x_j)-u^+(x_{j-1})|$ over finite partitions $a<x_0<\cdots<x_N<b$, equals $|Du|(I)$.
The proof below derives the interval statement directly from the distributional BV framework of \cite{Sim18}.
\end{theorem}
\begin{proof}
Subtract the cumulative function in Lemma \ref{lem:interval-measure-primitive}. The difference has zero distributional derivative, so Lemma \ref{lem:weak-gradient-zero} gives the constant $c$. The one-sided limits and jump identity follow from that lemma. For every $\varepsilon>0$, at most $|\mu|(I)/\varepsilon$ points can have $|\mu(\{x\})|\geq\varepsilon$; a countable union gives countability.

Each increment is $\mu((x_{j-1},x_j])$, so the variation is at most $|\mu|(I)$. For the reverse inequality, take $\psi\in C_c(I)$ with $|\psi|\leq1$. Choose endpoints outside its support and subdivide their compact interval into intervals of length at most $h$. Step functions constant on these half-open intervals, zero outside, approximate $\psi$ uniformly as $h\to0$. Their integrals against $\mu$ have absolute value at most the sum of the absolute interval masses, hence at most the pointwise variation. Passing to the limit gives $|\int\psi\,d\mu|$ bounded by that variation. The continuous-test characterization of total variation in Proposition \ref{prop:bv-vector-representation}, or the scalar case of its proof, gives $|\mu|(I)$ as the supremum of these integrals. This proves equality. The right-continuous representative is unique: two such representatives equal almost everywhere agree at every interior point by taking a sequence of equality points approaching from the right.
\end{proof}

\begin{proposition}[Absolute continuity and Lipschitz bounds]\label{prop:interval-bv-ac-lipschitz}
For $u\in BV(I)$, its regulated representative is absolutely continuous on $[a,b]$, using its endpoint traces, if and only if $Du=f\mathcal L^1$ for some $f\in L^1(I)$. It is Lipschitz with constant at most $L$ if and only if $|Du|\leq L\mathcal L^1$ as measures on $I$.
\end{proposition}
\begin{proof}
If $Du=f\mathcal L^1$, Theorem \ref{thm:interval-bv-representative} gives $u^+(x)=c+\int_a^xf(t)\,dt$. The absolute continuity of the integral and the one-dimensional fundamental theorem of calculus give the assertion. Conversely, that theorem for absolutely continuous functions identifies the distributional derivative with the integrable classical derivative.

The measure bound gives, for $x<y$, $|u^+(y)-u^+(x)|=|Du((x,y])|\leq L(y-x)$, and endpoint limits give the same bound on the closure. Conversely, a Lipschitz function is absolutely continuous with classical derivative bounded by $L$ almost everywhere, by the line differentiation result used in the Lipschitz foundations. Its derivative measure therefore has the stated domination. An absolutely continuous derivative measure here means absolute continuity relative to Lebesgue measure; it is distinct from merely having no atoms.
\end{proof}

\begin{theorem}[Finite-perimeter sets in an interval]\label{thm:interval-finite-perimeter-sets}
A measurable $E\subset I$ has finite perimeter in $I$ if and only if it agrees almost everywhere with a finite union of intervals relative to $I$. Choose the intervals with disjoint interiors and with no artificial common endpoint separating two occupied intervals. Then
\[
 D\chi_E=\sum_{\text{interior left endpoints }\ell}\delta_\ell
 -\sum_{\text{interior right endpoints }r}\delta_r,
\]
and $P(E;I)$ is the number of those interior endpoints. Endpoints at $a$ or $b$ make no contribution to relative perimeter.
This is a direct one-dimensional consequence of the BV definition of perimeter in \cite{Sim18}; the proof is included below.
\end{theorem}
\begin{proof}
Assume finite perimeter. The right-continuous representative of $\chi_E$ takes only values zero and one: approach any point from the right through points where it agrees with the original characteristic function. The left limits also take only these values. Every nonzero jump consequently has size one in absolute value. There are finitely many such jumps, since their total absolute mass is at most $P(E;I)$.

Away from those finitely many points the representative is continuous: the two one-sided limits agree with its value. On each intervening interval, a continuous function with values in $\{0,1\}$ is constant. Thus $E$ is almost everywhere a finite union of occupied intervals. Its derivative there is exactly the displayed sum, by integration of a test derivative on each interval. The positive and negative endpoint atoms have distinct supports after the stated merging, so total variation counts them. Conversely, that same computation for any finite union produces a finite derivative measure and hence finite perimeter. In particular the empty set and the whole interval both have relative perimeter zero.
\end{proof}

\begin{proposition}[Zero extension and endpoint traces]\label{prop:interval-bv-zero-extension}
Let $u\in BV(a,b)$ and extend it by zero to $\widetilde u\in L^1(\mathbb R)$. Then
\[
 D\widetilde u=(Du)\lfloor(a,b)+u(a+)\delta_a-u(b-)\delta_b,
\]
where the first term means the derivative measure extended by zero. Consequently
\[
 |D\widetilde u|(\mathbb R)=|Du|((a,b))+|u(a+)|+|u(b-)|.
\]
\end{proposition}
\begin{proof}
Use the representation $u(x)=c+\mu((a,x])$ almost everywhere. For $\phi\in C_c^1(\mathbb R)$, Fubini as in Lemma \ref{lem:interval-measure-primitive}, now retaining the endpoint term, gives
\[
 -\int_a^bu\phi'
 =\int_{(a,b)}\phi\,d\mu+c\phi(a)-(c+\mu((a,b)))\phi(b).
\]
The two coefficients are the endpoint traces. The three measures live on pairwise disjoint Borel sets, so their total variations add. This is why zero extension is not generally variation preserving.
\end{proof}
